\section{Names and token references}

A \emph{token} is the QPU language's name for a logical storage location. For
the circuits in this guide, a token normally identifies one qubit wire.
Instructions refer to tokens instead of physical machine addresses. During
compilation, the system resolves each live token to an integer qubit index.
This indirection makes protocols readable and lets child processes have private
workspace without colliding with their parents.

\subsection{Token names}

A bare token is written as a whitespace-delimited name:
\begin{lstlisting}
Q
Control
CarryOut
scratch_2
0
\end{lstlisting}
Names are case-sensitive even though instruction names are not:
\code{Carry} and \code{carry} are distinct tokens. A name cannot contain
whitespace because whitespace separates language elements. Avoid names that
begin with a hyphen: inside a flag list, any hyphen-prefixed token is treated
as the beginning of another flag.

Tokens are allocated lazily. Mentioning a token in a gate, a state-setting
instruction, or \code{CREATETOKEN} gives it a wire when compilation needs one.
Consequently, source order affects the internal numeric wire index, but not the
meaning of a consistently named circuit.

\subsection{Direct references}

A direct reference is simply the token's name:
\begin{lstlisting}
SET Q 0p
H -I Q -O Q
MEASURE -I Q
RETURNVALS Q
\end{lstlisting}
Every occurrence resolves to the same token within that process invocation.
For a one-qubit gate, repeat the name after both \code{-I} and \code{-O}.
This does not copy the state. It says that Q is the instruction's incoming
wire and that the same Q wire is mutated in place.

\subsection{Parameter references with a dollar sign}

A leading dollar sign is the explicit parameter/reference spelling:
\begin{lstlisting}
PARAMS: Input:state
MAIN-PROCESS DollarReference
SET Local $Input
H -I $Input -O $Input
RETURNVALS Local
\end{lstlisting}
Address resolution removes the leading \code{$}. Thus \code{Input} and
\code{$Input} ultimately refer to the same declared parameter. The dollar
form is especially useful inside reusable child protocols because it visually
separates caller-supplied values from child-local workspace.

\code{SET Local \$Input} creates an alias: Local resolves to the same underlying
wire as Input. It does \emph{not} clone an unknown quantum state. Quantum states
cannot generally be copied because of the no-cloning theorem; aliasing avoids
claiming that a physical copy occurred.

\subsection{Cycle-qualified references with \code{:}}

A token may carry a suffix after a colon:
\begin{lstlisting}
Q:0
Q:1
Carry:next
\end{lstlisting}
The current system strips the colon and everything after it when mapping source
references to visual qubit wires. Therefore all three examples above whose
base is Q---\code{Q:0}, \code{Q:1}, and bare \code{Q}---identify the same
logical wire. The suffix is protocol notation for a stage, cycle, or bit
version; it does not create a second qubit.

\begin{lstlisting}
SET Q:0 0p
H -I Q:0 -O Q:1
MEASURE -I Q:2
\end{lstlisting}
This example applies H to one wire and then measures that same wire. The
changing suffix communicates progression to a human reader, while
\code{INCREASECYCLE} provides the actual compiler cycle boundary.

\textbf{Practical rule:} if two values must exist simultaneously, give them
different base names. Do not expect \code{Q:0} and \code{Q:1} to hold two
independent states.

\subsection{Aliases created by \code{SET}}

When the second operand of SET is not a built-in state constant, SET creates a
name alias:
\begin{lstlisting}
SET Working Input
SET Output Working
X -I Output -O Output
\end{lstlisting}
After resolution, Working, Output, and Input name one wire. Mutating any alias
mutates that wire for every alias. This behavior is useful for readable names
and process binding, but it is not assignment in the classical
``copy a value into a new variable'' sense.

Aliases are resolved within a process frame. PARAMS and parent/child bindings
take precedence over ordinary child-local scoping, ensuring that an operation
on a bound parameter affects the caller's intended wire.

\subsection{Numeric workspace references}

Bare numeric names such as \code{0}, \code{1}, and \code{2} are valid token
names and are common in reusable arithmetic children:
\begin{lstlisting}
SET 0:0 $A
SET 1:0 $B
SET 2:0 $Cin
\end{lstlisting}
When a child process uses a numeric workspace name, the compiler places it in
the parent invocation's workspace namespace. This lets nested adders reuse the
same compact source names without globally identifying every child invocation's
scratch wires.

Numeric token names are not number literals. \code{2} names a wire; it does not
mean the integer two or the basis state $\ket{2}$.

\subsection{Process scope and child-local references}

Each process invocation receives a unique scope. A local name such as
\code{Scratch} in two separate child calls therefore denotes two separately
resolved local references unless output or parameter binding deliberately joins
them. Conceptually, compilation turns:
\begin{lstlisting}
RUNCHILD Worker -I A -O First
RUNCHILD Worker -I B -O Second
\end{lstlisting}
into two independent Worker frames. This is similar to local variables in two
function calls. PARAMS map caller references into the frame, and RETURNVALS map
selected frame references back out.

